\documentclass[11pt]{article}
\usepackage{times}
\usepackage{fullpage}
\usepackage{url}
\usepackage{hyperref}
\usepackage{fancyhdr}
\usepackage{graphicx}
\usepackage{enumitem}
\usepackage{subcaption}
\usepackage{amsmath, amsfonts, amsthm, fouriernc}
\usepackage{IEEEtrantools}
\usepackage{parskip}

\pagestyle{fancy}
\addtolength{\headheight}{2em}
\addtolength{\headsep}{1.5em}
\lhead{\large{ME 2801}}

\usepackage{sectsty}
\sectionfont{\fontsize{12}{8}\selectfont}

\begin{document}

\begin{center}
  \Large{\bf{Lab Assignment: USV PID Control}}
\end{center}

This assignment accompanies the closed-loop PID control experiments on the USV.
You will design controllers analytically, implement them in ArduPilot, collect
step-response data, and compare measured performance against predictions.

\textbf{Deliverable.}  A brief lab report (PDF) with figures and MATLAB code
appended.

\bigskip\hrule\bigskip

% ---------------------------------------------------------------
\section*{1\quad Pre-lab: Controller Design}

\begin{enumerate}

\item Using the surge and yaw plant models identified in the previous lab,
  design a proportional controller for each axis that achieves the rise-time
  specification given in class.  Show your design calculations.

\item For the surge axis, design a PI controller that eliminates steady-state
  error while maintaining approximately the same rise time.  State your choice
  of $K_p$ and $K_i$.

\item Predict the closed-loop step-response metrics (rise time, settling time,
  percent overshoot, steady-state error) for each controller.  Use MATLAB
  \texttt{step()} and \texttt{stepinfo()} to verify your predictions.

\end{enumerate}

% ---------------------------------------------------------------
\section*{2\quad Experiment}

\begin{enumerate}[resume]

\item Configure the ArduPilot parameters for your designed controller gains.
  Record the parameter file for each trial.

\item Conduct step-response experiments for the surge and yaw axes.
  Collect ArduPilot \texttt{.BIN} logs for each trial.

\end{enumerate}

% ---------------------------------------------------------------
\section*{3\quad Post-lab: Data Analysis}

\begin{enumerate}[resume]

\item Process the \texttt{.BIN} logs using the provided pipeline
  (\texttt{bin2mat.py}, \texttt{proc\_openloop\_response.m}).
  Plot the measured step response for each trial.

\item Overlay the simulated (model-predicted) response on each measured
  response plot.  Compute the rise time, settling time, and steady-state
  error from the measured data.

\item Fill in the comparison table below (expand rows as needed):

  \begin{center}
  \begin{tabular}{lccc}
    \hline
    Metric & Predicted & Measured & Difference \\
    \hline
    Surge $T_r$ (s) & & & \\
    Surge $e_{ss}$ (m/s) & & & \\
    Yaw $T_r$ (s) & & & \\
    Yaw $e_{ss}$ (rad/s) & & & \\
    \hline
  \end{tabular}
  \end{center}

\end{enumerate}

% ---------------------------------------------------------------
\section*{4\quad Discussion}

\begin{enumerate}[resume]

\item Discuss the agreement (or disagreement) between predicted and measured
  performance.  What are the most likely sources of discrepancy?

\item How did the PI controller affect steady-state error and rise time compared
  to the proportional-only controller?  Is this consistent with your
  pre-lab prediction?

\end{enumerate}

\end{document}
